\documentclass[12pt]{article}
\usepackage[utf8]{inputenc}
\usepackage[T1]{fontenc}
\usepackage{hyperref}
\title{Sob Language Resources for K-12 Curriculum Development}
\author{Compiled by Claude Sonnet 5 (claude-sonnet-5) for Andr\'as Kornai}
\date{}
\begin{document}
\maketitle

\section{Introduction}
Sob (ISO 639-3 urw) is a Trans--New Guinea language of Madang Province, Papua New Guinea. It is also called Sop, Usino and Kari, though speakers do not accept Usino as the name of their language. About 2,300 speakers were reported in 2003. Do not confuse it with Sobei (ISO sob), an Austronesian language of Papua, Indonesia.

\section{Primary Educational Resources}
I found no primers, textbooks or school programmes for Sob.

\section{Linguistic and Reference Materials}
\begin{itemize}
\item Overview of the language: \url{https://en.wikipedia.org/wiki/Sop_language}
\item Ryan Harty (2021), a conference paper on Sob syllable profile and glide formation, listed in Glottolog: \url{https://glottolog.org/resource/reference/id/587460}
\end{itemize}

\section{Religious and Community Resources}
None found.

\section{Digital and Audio Resources}
\begin{itemize}
\item Rosetta Project entry for urw: \url{https://rosettapanglossia.longnow.org/index.php?title=Urw}
\end{itemize}

\section{Curriculum Development Context}
Very little documentation turned up, and the only academic work found is a phonology paper. A curriculum would first need community-led documentation, such as vocabulary, stories and an agreed orthography.
\end{document}